% book04.tex --- Book IV of Euclid's Elements: Inscribed and Circumscribed Figures.
%
% All 16 propositions encoded. Book IV builds entirely on Books I and III,
% culminating in the inscription of the regular pentadecagon (IV.16) via
% combination of the inscribed equilateral triangle (IV.2) and the regular
% pentagon (IV.11).
%
% Wording follows Heath (1908); proof sketches mirror his density.

\section{Book IV --- Inscribed and Circumscribed Figures}
\label{sec:book-IV}

\begin{claim}[Proposition IV.1: Fit a chord of given length in a circle]
\label{prop:IV.1}
Into a given circle to fit a straight line equal to a given straight
line which is not greater than the diameter of the circle.
\end{claim}
\begin{evidence}[Proof of IV.1]
\label{ev:IV.1}
Let $ABC$ be the given circle and $D$ the given segment, not greater
than the diameter $BC$ of the circle.  Apply I.3 to cut off from $BC$
a segment $BE$ equal to $D$.  With centre $B$ and radius $BE$ describe
a circle $EAF$ (Postulate 3); join $BA$ where $EAF$ meets the given
circle.  Then $BA = BE = D$ by Definition I.15, so $BA$ is the required
chord.
\dependson{IV.1}{I.3}
\dependson{IV.1}{post:3}
\dependson{IV.1}{def:I.15}
\end{evidence}

\begin{claim}[Proposition IV.2: Inscribe a triangle similar to a given triangle in a circle]
\label{prop:IV.2}
In a given circle to inscribe a triangle equiangular with a given
triangle.
\end{claim}
\begin{evidence}[Proof of IV.2]
\label{ev:IV.2}
Let $ABC$ be the given circle and $DEF$ the given triangle.  Draw the
tangent $GH$ at any point $A$ of the circle (III.16, III.17).  On $GH$
construct $\angle HAC = \angle DEF$ (I.23), and $\angle GAB = \angle
DFE$.  Join $BC$.  By III.32 (tangent--chord angle equals the
inscribed angle in the alternate segment), $\angle ACB = \angle HAB
= \angle DEF$ and $\angle ABC = \angle GAC = \angle DFE$.  By I.32 the
remaining angles agree; so $\triangle ABC$ is equiangular with
$\triangle DEF$.
\dependson{IV.2}{I.23}
\dependson{IV.2}{I.32}
\dependson{IV.2}{III.16}
\dependson{IV.2}{III.32}
\end{evidence}

\begin{claim}[Proposition IV.3: Circumscribe a triangle about a circle]
\label{prop:IV.3}
About a given circle to circumscribe a triangle equiangular with a
given triangle.
\end{claim}
\begin{evidence}[Proof of IV.3]
\label{ev:IV.3}
Let $ABC$ be the given circle with centre $K$ and $DEF$ the given
triangle.  Produce $EF$ both ways to $G$, $H$.  At the centre $K$
construct $\angle AKB = \angle DEG$ and $\angle AKC = \angle DFH$
(I.23).  At $A$, $B$, $C$ draw the tangents to the circle (III.16,
III.17); they bound a triangle.  Because each tangent is perpendicular
to the radius at the point of contact (III.18), the angles of the
constructed triangle are the supplements of the central angles $\angle
AKB$, $\angle AKC$, $\angle BKC$, hence equal to the angles of
$\triangle DEF$ by I.13 and the construction.
\dependson{IV.3}{I.13}
\dependson{IV.3}{I.23}
\dependson{IV.3}{I.32}
\dependson{IV.3}{III.16}
\dependson{IV.3}{III.18}
\end{evidence}

\begin{claim}[Proposition IV.4: Inscribe a circle in a triangle]
\label{prop:IV.4}
In a given triangle to inscribe a circle.
\end{claim}
\begin{evidence}[Proof of IV.4]
\label{ev:IV.4}
Let $ABC$ be the given triangle.  Bisect $\angle ABC$ and $\angle BCA$
by $BD$ and $CD$ (I.9), meeting at $D$.  Drop perpendiculars $DE$,
$DF$, $DG$ from $D$ to $AB$, $BC$, $CA$ (I.12).  In the pairs of
triangles formed at $D$, the two angles and a common side give
congruence (I.26), whence $DE = DF = DG$.  The circle with centre $D$
and radius $DE$ touches each side (since the perpendiculars at the
feet make the sides tangents by III.16) and is inscribed in
$\triangle ABC$.
\dependson{IV.4}{I.9}
\dependson{IV.4}{I.12}
\dependson{IV.4}{I.26}
\dependson{IV.4}{III.16}
\end{evidence}

\begin{claim}[Proposition IV.5: Circumscribe a circle about a triangle]
\label{prop:IV.5}
About a given triangle to circumscribe a circle.
\end{claim}
\begin{evidence}[Proof of IV.5]
\label{ev:IV.5}
Let $ABC$ be the given triangle.  Bisect $AB$ at $D$ and $AC$ at $E$
(I.10).  From $D$ and $E$ draw perpendiculars to $AB$ and $AC$
respectively (I.11), meeting at $F$.  Join $FA$, $FB$, $FC$.  By I.4
applied to the two right triangles at $D$, $FA = FB$; similarly at
$E$, $FA = FC$.  Thus $FA = FB = FC$, and the circle with centre $F$
and radius $FA$ passes through all three vertices.
\dependson{IV.5}{I.4}
\dependson{IV.5}{I.10}
\dependson{IV.5}{I.11}
\dependson{IV.5}{def:I.15}
\end{evidence}

\begin{claim}[Proposition IV.6: Inscribe a square in a circle]
\label{prop:IV.6}
In a given circle to inscribe a square.
\end{claim}
\begin{evidence}[Proof of IV.6]
\label{ev:IV.6}
Let $ABCD$ be the given circle with centre $E$.  Draw two diameters
$AC$ and $BD$ at right angles (I.11).  Join $AB$, $BC$, $CD$, $DA$.
The four right triangles at $E$ are congruent by I.4 (two radii and
the common right angle), so $AB = BC = CD = DA$.  The inscribed
angles standing on the diameters are right (III.31), so all four
angles of $ABCD$ are right.  Therefore $ABCD$ is a square.
\dependson{IV.6}{I.4}
\dependson{IV.6}{I.11}
\dependson{IV.6}{III.31}
\dependson{IV.6}{def:I.15}
\end{evidence}

\begin{claim}[Proposition IV.7: Circumscribe a square about a circle]
\label{prop:IV.7}
About a given circle to circumscribe a square.
\end{claim}
\begin{evidence}[Proof of IV.7]
\label{ev:IV.7}
Draw two perpendicular diameters of the given circle (I.11).  Through
each endpoint draw the tangent to the circle (III.16).  Each tangent
is perpendicular to its diameter (III.18); the four tangents thus form
a quadrilateral with all sides parallel and all angles right (I.28,
I.29).  Equal radii combined with I.34 force the sides equal, hence a
square circumscribes the circle.
\dependson{IV.7}{I.11}
\dependson{IV.7}{I.28}
\dependson{IV.7}{I.29}
\dependson{IV.7}{I.34}
\dependson{IV.7}{III.16}
\dependson{IV.7}{III.18}
\end{evidence}

\begin{claim}[Proposition IV.8: Inscribe a circle in a square]
\label{prop:IV.8}
In a given square to inscribe a circle.
\end{claim}
\begin{evidence}[Proof of IV.8]
\label{ev:IV.8}
Let $ABCD$ be the given square.  Bisect the sides at $E$, $F$, $G$,
$H$ (I.10).  Join $EG$ and $FH$, meeting at $K$.  By I.34 and the
construction, $KE = KF = KG = KH$.  Drop perpendiculars from $K$ to
each side; each is equal to $KE$.  Thus the circle with centre $K$
and radius $KE$ touches each side at its midpoint (III.16).
\dependson{IV.8}{I.10}
\dependson{IV.8}{I.34}
\dependson{IV.8}{III.16}
\end{evidence}

\begin{claim}[Proposition IV.9: Circumscribe a circle about a square]
\label{prop:IV.9}
About a given square to circumscribe a circle.
\end{claim}
\begin{evidence}[Proof of IV.9]
\label{ev:IV.9}
Join the diagonals $AC$ and $BD$ of the given square, intersecting
at $E$.  In the right triangles $\triangle ABE$ and $\triangle ADE$,
SSS (I.8) gives $\angle EAB = \angle EAD$, so $AE$ bisects the right
angle at $A$; similarly at every vertex.  The four triangles at $E$
are then isosceles with equal vertex angles (I.6), so $EA = EB = EC =
ED$.  The circle with centre $E$ and radius $EA$ passes through all
four vertices.
\dependson{IV.9}{I.6}
\dependson{IV.9}{I.8}
\dependson{IV.9}{def:I.15}
\end{evidence}

\begin{claim}[Proposition IV.10: Construct an isosceles triangle with 72--72--36 angles]
\label{prop:IV.10}
To construct an isosceles triangle having each of the angles at the
base double of the remaining one.
\end{claim}
\begin{evidence}[Proof of IV.10]
\label{ev:IV.10}
Take a straight line $AB$ and cut it at $C$ so that the rectangle on
$AB$, $BC$ equals the square on $AC$ (II.11, the golden section).
With centre $A$ and radius $AB$ describe a circle; in it apply chord
$BD$ equal to $AC$ (IV.1).  Join $AD$, $CD$.  Because $AB \cdot BC =
AC^2 = BD^2$, $BD$ is tangent to the circle through $A$, $C$, $D$
(III.37); by III.32 the tangent--chord angle equals the alternate
inscribed angle.  Tracking the resulting angle relations (with I.5
for the isosceles base angles) gives the required ratio.
\dependson{IV.10}{II.11}
\dependson{IV.10}{IV.1}
\dependson{IV.10}{IV.5}
\dependson{IV.10}{I.5}
\dependson{IV.10}{I.32}
\dependson{IV.10}{III.32}
\dependson{IV.10}{III.37}
\end{evidence}

\begin{claim}[Proposition IV.11: Inscribe a regular pentagon in a given circle]
\label{prop:IV.11}
In a given circle to inscribe an equilateral and equiangular pentagon.
\end{claim}
\begin{evidence}[Proof of IV.11]
\label{ev:IV.11}
Construct the 72--72--36 isosceles triangle $FGH$ by IV.10.  Inscribe
in the given circle a triangle $ACD$ equiangular with $FGH$ (IV.2).
Bisect the base-angles of $ACD$ by IV.10's construction propagated
into the circle, yielding two more division points $B$, $E$.  The
five arcs are equal (III.26), so the five chords $AB$, $BC$, $CD$,
$DE$, $EA$ are equal (III.29), and the inscribed angles standing on
equal arcs are equal (III.27): the pentagon is regular.
\dependson{IV.11}{I.9}
\dependson{IV.11}{IV.2}
\dependson{IV.11}{IV.10}
\dependson{IV.11}{III.26}
\dependson{IV.11}{III.27}
\dependson{IV.11}{III.29}
\end{evidence}

\begin{claim}[Proposition IV.12: Circumscribe a regular pentagon about a circle]
\label{prop:IV.12}
About a given circle to circumscribe an equilateral and equiangular
pentagon.
\end{claim}
\begin{evidence}[Proof of IV.12]
\label{ev:IV.12}
Inscribe a regular pentagon in the given circle by IV.11.  At each
vertex draw the tangent (III.16); the five tangents bound the
circumscribed pentagon.  Each tangent is perpendicular to its radius
(III.18), and by I.4 the right triangles formed at adjacent vertices
are congruent, so the circumscribed pentagon has equal sides and
equal angles.
\dependson{IV.12}{IV.11}
\dependson{IV.12}{I.4}
\dependson{IV.12}{III.16}
\dependson{IV.12}{III.18}
\end{evidence}

\begin{claim}[Proposition IV.13: Inscribe a circle in a regular pentagon]
\label{prop:IV.13}
In a given pentagon, which is equilateral and equiangular, to inscribe
a circle.
\end{claim}
\begin{evidence}[Proof of IV.13]
\label{ev:IV.13}
Bisect two adjacent interior angles of the pentagon (I.9); their
bisectors meet at a point $F$.  Drop perpendiculars from $F$ to each
side (I.12); by I.4 these perpendiculars are equal.  The circle on
$F$ with that common radius touches every side (III.16).
\dependson{IV.13}{I.4}
\dependson{IV.13}{I.9}
\dependson{IV.13}{I.12}
\dependson{IV.13}{III.16}
\end{evidence}

\begin{claim}[Proposition IV.14: Circumscribe a circle about a regular pentagon]
\label{prop:IV.14}
About a given pentagon, which is equilateral and equiangular, to
circumscribe a circle.
\end{claim}
\begin{evidence}[Proof of IV.14]
\label{ev:IV.14}
Take the same point $F$ as in IV.13 (intersection of two
angle-bisectors).  Join $F$ to each vertex; by I.4 the resulting
triangles are congruent (equal sides, common bisected angles), so the
five distances from $F$ to the vertices are equal.  Draw the circle
on $F$ with that radius (Definition I.15).
\dependson{IV.14}{IV.13}
\dependson{IV.14}{I.4}
\dependson{IV.14}{def:I.15}
\end{evidence}

\begin{claim}[Proposition IV.15: Inscribe a regular hexagon in a circle]
\label{prop:IV.15}
In a given circle to inscribe an equilateral and equiangular hexagon.
\end{claim}
\begin{evidence}[Proof of IV.15]
\label{ev:IV.15}
Let $G$ be the centre and $AD$ a diameter of the given circle.  With
centre $D$ and radius $DG$ describe a circle meeting the given circle
at $C$ and $E$ (Postulate 3).  Join $GC$, $GE$.  The triangle $GDC$
has $GD = DC = CG$ (radii of equal circles), so it is equilateral,
and $\angle DGC = 60^\circ$ (I.32).  Stepping this $60^\circ$ chord
around the circle six times produces the regular hexagon.
\dependson{IV.15}{I.32}
\dependson{IV.15}{post:3}
\dependson{IV.15}{def:I.15}
\end{evidence}

\begin{claim}[Proposition IV.16: Inscribe a regular 15-gon in a circle]
\label{prop:IV.16}
In a given circle to inscribe a fifteen-angled figure which shall be
both equilateral and equiangular.
\end{claim}
\begin{evidence}[Proof of IV.16]
\label{ev:IV.16}
Inscribe a regular pentagon (IV.11) and a regular equilateral
triangle (IV.2) in the circle, sharing a common vertex $A$.  The arc
from $A$ to the next pentagon-vertex is $\tfrac{1}{5}$ of the circle;
the arc from $A$ to the next triangle-vertex is $\tfrac{1}{3}$.  The
difference is $\tfrac{1}{3} - \tfrac{1}{5} = \tfrac{2}{15}$ of the
circle.  Bisect that arc (III.30); each half is $\tfrac{1}{15}$ of
the circle, and stepping that chord fifteen times around gives the
regular pentadecagon.
\dependson{IV.16}{IV.2}
\dependson{IV.16}{IV.11}
\dependson{IV.16}{III.30}
\end{evidence}
